\subsection{Justificativa}

\par Os processos industriais geram grande volume de dados, apresentados ao operador por um software de supervisão, por meio de uma interface homem-máquina (\sigla{ihm}), que converte leituras de campo em telas, tendências e alarmes \cite{ZANGHI2019_SCADA_CONCEITOS}. Com computadores mais capazes, protocolos melhores e a integração com a Internet, consolidou-se a supervisão remota. Em setores como saneamento, transporte e energia, com instalações espalhadas, ganham peso a telemetria e o controle à distância. Uma arquitetura típica de \sigla{scada} integra os equipamentos de campo, como CLPs, aos módulos de monitoramento e controle e centraliza a visualização para o operador \cite{AHMED2011_EFFICIENT_WEB_SCADA}.

\par Os supervisórios convencionais costumam exigir licenças proprietárias, hardware dedicado e infraestrutura local, o que os mantém fora do alcance de pequenas e médias empresas \cite{PHUYAL2020_COST_EFFICIENT_SCADA}. Oferecer a supervisão pelo navegador reduz esse investimento e amplia o monitoramento a instalações dispersas, com menos deslocamentos e resposta mais ágil a ocorrências \cite{AHMED2011_EFFICIENT_WEB_SCADA}.

\par A computação em nuvem oferece armazenamento, processamento, bancos de dados e aplicações como serviços remotos \cite{PEDROSA2011_COMPUTACAO_NUVEM}, com capacidade sob demanda, o que reduz o investimento local \cite{HENRIQUES2010_CLOUD_COMPUTING_TCC}. A conectividade, em especial sem fio, permite que dispositivos e sistemas ciber-físicos troquem dados com a nuvem \cite{STANKOVIC2014_IOT_RESEARCH_DIRECTIONS,MONOSTORI2016_CPS_MANUFACTURING}.

\par Neste trabalho, a nuvem é uma possibilidade de aplicação, e não a forma de implantação adotada. O sistema roda em contêineres num servidor, o que dispensa contratos de hospedagem e mantém baixo o custo, requisito central para pequenas e médias empresas. Como cada serviço é um contêiner independente, a solução poderia ir a um provedor de nuvem sem alteração de código \cite{NETTO2017_COORDENACAO_CONTAINERS_KUBERNETES,NETTO2016_STATE_MACHINE_REPLICATION_KUBERNETES}, e o acesso pelo navegador dispensa instalar programas nos computadores dos usuários.
